\enteteExam{NFE106 -- Administration de bases de données}%
{Examen 13 septembre 2024 -- session 2 -- 3h}%
{Examen en présentiel -- Calculatrices autorisées -- Documents autorisés: 5 pages recto/verso de notes personnelles}%


Le sch\'ema de base de donn\'ees suivant sera utilis\'e pour l'ensemble du
sujet. Il permet de gérer une base de données d'une base gérant une archive de publications scientifiques (20 ans de publications).

\begin{itemize}
\item {\sc Publication} (\underline{IdPub}, Titre, Date, IdJournal), avec $3,15 \times 10^7$ enregistrements
\item \textsc{Journal}  (\underline{IdJournal}, Titre, Acronyme, Rang), avec  26 400 enregistrements
\item  \textsc{Citation}  (\underline{IdPub, IdPubCitée}), avec $5,53 \times 10^8$ enregistrements
\item  \textsc{Auteur}  (\underline{IdPub, position}, nomAuteur, prenomAuteur)  $1,55 \times 10^8$ enregistrements
\end{itemize}

\smallskip

On suppose que les paramètres de stockage sont les suivants\,:
\begin{itemize}
  \item 1 octet\,: position
	\item 3 octets : Date
	\item 4 octets (entiers) : id, position (max 40)
	\item 10 octets (chaînes courtes) : Rang (A*, A, B, C, Inconnu), Acronyme
	\item 50 octets : Titre, Nom, Prenom
\end{itemize}


Autres statistiques\,:
\begin{itemize}
	\item 2M d'auteurs\,;
	\item l'acronyme est unique\,;
	\item dans la table \textsc{Auteur} le triplet (idPub, nomAuteur, prenomAuteur) est unique.
\end{itemize}

Les adresses physiques sur disque occupent 10 octets, et un bloc 8\,192 octets.

\section{Stockage et indexation (8 points)}


\begin{enumerate}
	\item {\tt (2~points)} Présenter sous forme de tableau la taille (en nombre de
    blocs) prise par chaque table, en supposant qu'on laisse un espace
    libre de 10\% et qu'il n'y a pas d'entête\,;
\begin{solution}
	L'espacece utile dans un bloc est de 7372 octets.\\
	\begin{tabular}{rrrrr}
\textbf{Table} & {\bf Taille tuples} 	& {\bf Nb tuples par bloc} & \textbf{Nb tuples} & {\bf Nb blocs}\\
\hline
Publication&	68	&	105	&	31 500 000	&	300 000	\\
Journal&	81	&	88	&	26 400	&	300	\\
Citation&	13	&	553	&	553 000 000	&	1 000 000	\\
Auteur&	115	&	62	&	155 000 000	&	2 500 000	\\
\end{tabular}
\end{solution}

	\item {\tt (2~points)} La sélectivité d'un attribut est le nombre d'enregistrements
     (en moyenne)
     ramenée par une recherche sur une table pour une valeur de
     cet attribut.  Donner la sélectivité  de\,:
	\begin{itemize}
		\item \textsc{Publication}.IdJournal\,;
		\item \textsc{Journal}.Acronyme\,;
		\item \textsc{Auteur}.IdPub\,;
		\item \textsc{Auteur}.(NomAuteur, prenomAuteur)\,;
	\end{itemize}
\begin{solution}
\begin{itemize}
	\item 1
	\item 1
	\item 1
	\item $155 10^6 /(2\times 10^6) = 77.5$
\end{itemize}
\end{solution}

	\item {\tt (2~points)} Donner la hauteur d'un index dense sur la clé primaire de ``\textsc{Auteur}''.
    Déterminez d'abord la taille d'une entrée de l'index puis
    le nombre d'entrées par bloc.  Vous en déduirez la hauteur par
    calcul ``à la main'', niveau par niveau, ou en appliquant le logarithme
    approprié.

\begin{solution}
\begin{itemize}
	\item clé : 3+(1+4) + (1+1) + 10
	\item ordre : 180
	\item $log_{180}(1,55 \times 10^8) = 4$
\end{itemize}
\end{solution}

	\item {\tt (2~points)} Donner l'arbre B  d'ordre 2 sur \textsc{Auteur}.(IdPub, position)
  pour les valeurs suivantes\,:\\

  \begin{verbatim}
      (100-2), (350-1), (240-1), (200-1), (280-5), (100-1),
      (270-1), (200-2), (200-3),(350-2), (240-3), (220-1)
  \end{verbatim}

  \begin{solution}
\begin{figure}
	\centering
%	\includegraphics[width=0.7\linewidth]{figs/2018-06_btree}
	\caption{}
	\label{fig:2018-06btree}
\end{figure}
\end{solution}

\end{enumerate}

\section{Optimisation (7 points)}

\begin{enumerate}
	\item {\tt (3~points)} Soit la requête suivante\,:
\begin{minted}{sql}
SELECT idJournal FROM Journal WHERE Titre like '%database%' AND Rang = 'A*';
\end{minted}
	Donner le coût d'évaluation (nombre de lectures de blocs au pire des cas)
   pour chacun des scénarios  suivants :
	\begin{enumerate}
		\item Il n'existe qu'un arbre  B (non dense) sur le titre
		\item  Il n'existe qu'un arbre  B (non dense) sur le rang, avec la distribution suivante\,:\\
		\begin{tabular}{|c|c|c|c|c|}
			\hline
			A*& A& B& C& Inconnu\\
			\hline
			7,644\%& 15\% & 20\% & 25\% & 32,356\%\\
			\hline
		\end{tabular}
    \item Il existe qu'un arbre  B sur le rang, mais cette fois il est dense.
     Expliquez ce que cela implique, et donnez le nouveau calcul.
		\item Il existe une table de hachage de taille 5 sur
    l'attribut Rang (1 rang = 1 partition). Donner le coût précis pour cette requête.
	\end{enumerate}
	\begin{solution}
		\begin{enumerate}
			\item L'arbre  ne sert à rien, il faut un parcours séquentiel\,: 300 I/O
			\item On peut utiliser l'arbre  sur la rang, mais il est très peu
      efficace comme le montre le calcul suivant.
      Il existe environ $26\, 400*7,644/100$ journaux de rang A*,
      donc, au pire, à peu près autant de lectures de blocs (2000).

      \item Si l'arbre est dense, on trouve les 2,000 journaux de rang A*
         séquentiellement, soit $2000 / 88 \approx 23$ blocs
			\item Avec une table de hachage le résultat est à peu
         près le même (au parcours d'arbre près)
         puisque tous les journaux de rang A* sont dans le même fragment.
		\end{enumerate}
	\end{solution}

	\item \texttt{(2~points)} Soit le plan EXPLAIN suivant.
  Expliquez-le et donnez la requête SQL, sachant que l'on veut grouper et compter le résultat par IdJournal\,:
\begin{verbatim}
0   SELECT STATEMENT
1     HASH GROUP BY
2*       HASH JOIN
3*          PARTITION HASH SINGLE
4*              TABLE ACCESS FULL             Journal
5           TABLE ACCESS BY INDEX ROWIDS      Publication
6*             INDEX RANGE SCAN               IDX_Pub
(2) J.IdJournal = P.IdJournal
(3) Rang = 'A*'
(4) Titre like '%Database%'
(6) Date like '2018-%'
\end{verbatim}

\begin{solution}
\begin{minted}{sql}
SELECT IdJournal, COUNT(*)
FROM Journal J, Publication P
WHERE J.IdJournal = P.IdJournal
   AND J.Rang = 'A*' AND J.titre like '%Database%%'
   AND P.Date like '2018-%'
GROUP BY idjournal
\end{minted}
\end{solution}

	\item \texttt{(2~points)} Proposer une amélioration
  du plan d'exécution précédent en supposant qu'il existe seulement un index
  sur chaque clé primaire.
\begin{solution}
On effectue une  Boucle imbriquée avec index  sur (Id Journal, date)
en parcourant les publications et en accédant aux journaux
par l'index.
\end{solution}
\end{enumerate}


\section{Concurrence (5 points)}

Soit l'ex\'ecution concurrente suivante de trois transactions dans un SGBD.

$$H: r_1[x] r_2[x] r_1[y] w_2[y] r_3[x] w_1[x] r_3[y] C_1 C_2 C_3$$

\begin{enumerate}
   \item Existe-t-il une ``lecture sale''~? Pourquoi faut-il \'eviter
   ce type de lecture (expliquer les effets possibles)\,? (1~point)

    \item Cette exécution est-elle sérialisable (1~point)\,?

     \item Quelle est l'ex\'ecution obtenue par verrouillage \`a deux phases
   \`a partir de $H$?  (2~points)

   \item Quelle transaction pourrait effectuer une lecture avec la clause
     \texttt{for update} et quel en serait l'effet\,?
\end{enumerate}

\begin{solution}
\begin{enumerate}
   \item Oui, $r_3[y]$ est une lecture sale (on lit la version non validée
    de $T_2$)\,: si $T_2$ choisit d'annuler, $T_3$ s'appuie sur une
      valeur qui n'existe plus.

    \item Il y a un cycle dans les conflits entre $T_1$ et $T_2$.

     \item Interlocage entre $T_1$ et $T_2$.

   \item $T_1$ pourrait utiliser cette clause sur la lecture de $x$:
   $T_2$ serait alors mise en attente et il n'y aurait plus d'interlocage.
\end{enumerate}
\end{solution}
